\documentclass[12pt, a4paper]{article} %pour définir le type de document, la taille de la police et le format du papier

%Préambule : pour importer les packages, définir les commandes, les environnements, les styles, etc. qui seront utilisés dans le document

\usepackage[utf8]{inputenc} %pour pouvoir utiliser les caractères accentués directement dans le code source du document
\usepackage[T1]{fontenc} %pour que les caractères accentués soient correctement affichés dans le document final
\usepackage[french]{babel} %pour que les règles de typographie française soient appliquées dans le document (espaces insécables, guillemets français, etc.)
%\usepackage{hyperref} %pour créer des liens hypertexte
%\usepackage{graphicx} %pour inclure des images
%\usepackage{grafcet} %pour créer des diagrammes GRAFCET


%titre
\title{Hello World} 
\author{Pierre-Marie \textsc{Riboulet}} %pour définir l'auteur du document, le \textsc{} pour mettre le nom de famille en petites majuscules
%\date{} %pour ne pas afficher la date
%\date{02/02/2025} %pour afficher une date fixe


\begin{document}
\maketitle %affiche le titre, l'auteur et la date determiner si dessus


%\tableofcontents  %affiche la table des matières concu automatiquement

  %texte

   % \href{https://tiaportal_formation_automatisme.solutions-industrielles.com/}{lien vers le site de formation} %pour faire un lien vers un site internet copier le lien et changer les slash en antislash


   % \ref{Hello World} %pour faire référence à une section
%différents niveaux de section de textes :
%Partie
%(en class report ont a les chaptire) 
%section
%subsection
%subsubsection
%paragraph
%subparagraph

%\section*{Hello World} %\label{sec:patate} %le * pour ne pas numéroter la section, et le \label pour pouvoir faire référence à cette section plus tard dans le document, ou plutot a ce moment la compilée 2 fois.


%\emph{Hello World} \footnote{Hello world} %pour faire une note de bas de page, documenter une expression ou un mot
%\textbf{Gras, gros, fort} %pour faire du gras
%\textit{italique, penché} %pour faire de l'italique
%{\bf \it 
%}%pour faire du gras et de l'italique en même temps sur la section entre les accolades

    
%\bfseries %pour mettre en gras
%\itshape %pour mettre en gras et italique

%\normalfont %pour revenir à la police normale après avoir mis en gras et italique

% ou normalshape pour suprimer l'italique en gardant le gras
\end{document}

zone libre non pris en compte dans le document. ne sert pas à grand chose dans ce cas, mais peut être utile pour faire des notes ou des rappels dans le code source d'un autre document plus complexe.
pour compiler: pdflatex nonducument.tex
les sot de ligne et espace/tabulation sont ignorés dans le document, sauf si on utilise des commandes spécifiques pour les afficher (par exemple \newline pour un saut de ligne,
 ou \hspace{1cm} pour un espace horizontal de 1 cm).
commande: \nom de la commande[argulent suplémentaire] {argument obligatoire}
